%% ============================================================
%% GBC073 — Aula 5: Treinando redes II
%% Convergência, superfícies de perda e taxa de aprendizado
%% Baseada na Lecture 5 do CMU 11-785 (Convergence, Learning
%% Rates, and Gradient Descent), modernizada para PyTorch.
%%
%% Compilar:  latexmk -xelatex aula05.tex
%% ============================================================
\documentclass[aspectratio=169, 11pt]{beamer}

\makeatletter
\def\input@path{{./}{../}}
\makeatother

\usetheme{InteligenciaComputacional}   % ANTES do babel: fontspec vem primeiro
\usepackage{babel}
\babelprovide[import, main]{brazilian}

% ---- código Python/PyTorch (usa a paleta do tema) ----------
\usepackage{listings}
\lstset{
  language=Python,
  basicstyle=\ttfamily\scriptsize,
  keywordstyle=\color{icAzul}\bfseries,
  commentstyle=\color{icCinza}\itshape,
  stringstyle=\color{icVerde},
  emph={backward,step,zero_grad,SGD,Adam,AdamW,LambdaLR,parameters,no_grad},
  emphstyle=\color{icMagenta},
  showstringspaces=false,
  columns=fullflexible,
  keepspaces=true,
}

\title[Aula 5 — Convergência e taxa de aprendizado]%
      {Aula 5: Treinando redes II —\\ convergência, superfícies de perda e taxa de aprendizado}
\subtitle[GBC073 — Inteligência Computacional]%
         {Descida de gradiente \textbullet\ Número de condição \textbullet\ Momento}
\author[Prof.\ Marcelo Albertini]{Prof.\ Marcelo Keese Albertini}
\institute{GBC073 — Inteligência Computacional \textbullet\ Universidade Federal de Uberlândia}
\date{2026}

\begin{document}

% ------------------------------------------------------------ capa
\begin{frame}[plain, noframenumbering]
  \titlepage
\end{frame}

% ------------------------------------------------------------ onde estamos
\begin{frame}{Onde estamos}
  \framesubtitle{Aula 4: a retropropagação nos dá $\deriv{\erro}{\vw}$ para qualquer rede. E agora?}

  \begin{itemize}\setlength{\itemsep}{0.4ex}
    \item Já sabemos \textbf{representar}: MLPs aproximam qualquer função contínua
    \item Já sabemos \textbf{derivar}: retropropagação calcula \(\deriv{\erro}{\vw}\)
          em tempo linear no tamanho da rede
    \item A receita é iterar
          \(\vw \leftarrow \vw - \taxa \deriv{\erro}{\vw}\) \dots\ mas:
  \end{itemize}

  \begin{ideiabox}
    Ter o gradiente \emph{não} garante chegar a um bom mínimo, nem chegar rápido.
    A pergunta desta aula é: \textbf{quando e com que velocidade a descida de
    gradiente converge — e o que fazer quando ela não converge?}
  \end{ideiabox}

  \begin{itemize}\setlength{\itemsep}{0.2ex}
    \item Hoje: superfície de perda, análise de convergência, taxa de aprendizado,
          \emph{schedules} e momento
    \item Próxima aula: SGD, minilotes e otimizadores adaptativos (Adam, AdamW)
  \end{itemize}
\end{frame}

% ============================================================
\section{A superfície de perda}

\begin{frame}{O objeto que estamos descendo}
  \begin{defbox}{Superfície de perda}
    Dado um conjunto de dados \(\D = \{(\vx_i, \vy_i)\}_{i=1}^{N}\), a superfície
    de perda é a função
    \(\erro(\vw) = \frac{1}{N}\sum_{i=1}^{N} \ell\bigl(f(\vx_i; \vw),\, \vy_i\bigr)\),
    vista como função \emph{dos pesos} \(\vw \in \R^d\). Treinar é resolver
    \(\argmin_{\vw} \erro(\vw)\).
  \end{defbox}

  \begin{itemize}\setlength{\itemsep}{0.2ex}
    \item A entrada da rede é \(\vx\); a variável da otimização é \(\vw\),
          com \(d\) indo de milhões a \emph{trilhões}
    \item Não dá para visualizar — mas dá para \textbf{analisar}: a geometria
          dessa superfície decide velocidade, estabilidade e qualidade do treino
  \end{itemize}

  \begin{alertabox}
    Para redes com camadas ocultas, \(\erro(\vw)\) \alert{não é convexa}: há
    simetrias (permutar neurônios não muda a função), platôs, selas e vales
    curvos. Toda a teoria clássica de otimização convexa vale apenas
    \emph{localmente}.
  \end{alertabox}
\end{frame}

% ------------------------------------------------------------
\begin{frame}{Convexo vs.\ não convexo}
  \begin{columns}[T]
    \begin{column}{0.46\textwidth}
      \begin{center}
        {\color{icCinza}\bfseries\small Convexa: um único mínimo}\\[0.4em]
        \scalebox{0.8}{%
        \begin{tikzpicture}
          \draw[->, icCinza] (-1.5,0) -- (1.5,0) node[below] {\scriptsize $w$};
          \draw[->, icCinza] (0,-0.2) -- (0,1.9) node[left] {\scriptsize $\erro$};
          \draw[very thick, icAzul] plot[domain=-1.25:1.25, samples=40] (\x, {\x*\x});
          \fill[icVerde] (0,0) circle (2pt);
        \end{tikzpicture}}
      \end{center}
    \end{column}
    \begin{column}{0.5\textwidth}
      \begin{center}
        {\color{icCinza}\bfseries\small Não convexa: mínimos, platôs, selas}\\[0.4em]
        \scalebox{0.8}{%
        \begin{tikzpicture}
          \draw[->, icCinza] (-2.1,0) -- (2.3,0) node[below] {\scriptsize $w$};
          \draw[->, icCinza] (-2,-0.2) -- (-2,1.9);
          \draw[very thick, icMagenta]
            plot[domain=-1.9:2.1, samples=80]
            (\x, {0.9 + 0.5*sin(2.2*\x r) - 0.35*\x + 0.12*\x*\x});
          \fill[icDourado] (-1.0,{0.9 + 0.5*sin(-2.2 r) + 0.35 + 0.12}) circle (2pt);
          \fill[icVerde]   ( 1.62,{0.9 + 0.5*sin(3.564 r) - 0.567 + 0.315}) circle (2pt);
        \end{tikzpicture}}
      \end{center}
    \end{column}
  \end{columns}

  \medskip
  \begin{ideiabox}
    Visão moderna (Dauphin et al., 2014; Choromanska et al., 2015): em altíssima
    dimensão, o vilão típico \emph{não} é o mínimo local ruim — são os
    \textbf{pontos de sela} e os \textbf{platôs}, onde o gradiente quase some.
    Em redes largas, quase todos os mínimos locais têm perda parecida com a do
    mínimo global. Por isso a descida de gradiente funciona na prática.
  \end{ideiabox}
\end{frame}

% ============================================================
\section{Convergência do gradiente descendente}

\begin{frame}{O laboratório perfeito: a quadrática em 1D}
  \framesubtitle{Perto de um mínimo, toda superfície suave é aproximadamente quadrática (Taylor)}

  Tome \(\erro(w) = \tfrac{1}{2} a\, w^2\) com \(a > 0\)
  (curvatura \(\erro''(w) = a\)). A atualização
  \(w \leftarrow w - \taxa\, \erro'(w)\) vira:
  \[
    w_{t+1} \;=\; w_t - \taxa\, a\, w_t
    \;=\; \termo{(1 - \taxa a)}{fator de contração}\, w_t
    \qquad\Longrightarrow\qquad
    w_t = (1-\taxa a)^{t}\, w_0
  \]

  \begin{teobox}{Convergência na quadrática}
    A descida de gradiente converge se, e somente se, \(\abs{1 - \taxa a} < 1\),
    isto é, \(\mdestaque{icMagenta}{0 < \taxa < 2/a}\). Com
    \(\taxa = 1/a\) ela converge em \emph{um único passo}.
  \end{teobox}

  \begin{itemize}\setlength{\itemsep}{0.2ex}
    \item \(\taxa < 1/a\): converge de forma monotônica (mesmo lado do vale)
    \item \(1/a < \taxa < 2/a\): converge \emph{oscilando} de um lado para o outro
    \item \(\taxa > 2/a\): \alert{diverge geometricamente} — a perda explode
  \end{itemize}
\end{frame}

% ------------------------------------------------------------
\begin{frame}{Os três regimes, desenhados}
  \begin{columns}[T]
    \begin{column}{0.32\textwidth}
      \begin{center}
        {\color{icVerde}\bfseries\small $\taxa a = 0{,}5$: monotônico}\\[0.3em]
        \scalebox{0.72}{%
        \begin{tikzpicture}
          \draw[icCinza!50] plot[domain=-1.3:1.3, samples=40] (\x, {\x*\x});
          \draw[->, thick, icVerde] (1.0,1.0) -- (0.5,0.25);
          \draw[->, thick, icVerde] (0.5,0.25) -- (0.25,0.0625);
          \draw[->, thick, icVerde] (0.25,0.0625) -- (0.125,0.0156);
          \fill[icVerde] (1.0,1.0) circle (2pt);
        \end{tikzpicture}}
      \end{center}
    \end{column}
    \begin{column}{0.32\textwidth}
      \begin{center}
        {\color{icDourado}\bfseries\small $\taxa a = 1{,}7$: oscila e converge}\\[0.3em]
        \scalebox{0.72}{%
        \begin{tikzpicture}
          \draw[icCinza!50] plot[domain=-1.3:1.3, samples=40] (\x, {\x*\x});
          \draw[->, thick, icDourado] (1.0,1.0) -- (-0.7,0.49);
          \draw[->, thick, icDourado] (-0.7,0.49) -- (0.49,0.24);
          \draw[->, thick, icDourado] (0.49,0.24) -- (-0.34,0.117);
          \fill[icDourado] (1.0,1.0) circle (2pt);
        \end{tikzpicture}}
      \end{center}
    \end{column}
    \begin{column}{0.32\textwidth}
      \begin{center}
        {\color{icCarmim}\bfseries\small $\taxa a = 2{,}4$: diverge}\\[0.3em]
        \scalebox{0.72}{%
        \begin{tikzpicture}
          \draw[icCinza!50] plot[domain=-1.3:1.3, samples=40] (\x, {\x*\x});
          \draw[->, thick, icCarmim] (0.3,0.09) -- (-0.42,0.176);
          \draw[->, thick, icCarmim] (-0.42,0.176) -- (0.59,0.348);
          \draw[->, thick, icCarmim] (0.59,0.348) -- (-0.82,0.672);
          \draw[->, thick, icCarmim] (-0.82,0.672) -- (1.15,1.32);
          \fill[icCarmim] (0.3,0.09) circle (2pt);
        \end{tikzpicture}}
      \end{center}
    \end{column}
  \end{columns}

  \medskip
  \begin{ideiabox}
    A taxa de aprendizado ótima é o \emph{inverso da curvatura}:
    \(\taxa^\ast = 1/\erro''\). Vales muito curvos pedem passos pequenos; vales
    rasos pedem passos grandes. Guarde essa frase — ela explica o método de
    Newton, o RMSProp e o Adam.
  \end{ideiabox}
\end{frame}

% ------------------------------------------------------------
\begin{frame}{Muitas dimensões: autovalores e número de condição}
  Perto do mínimo, \(\erro(\vw) \approx \tfrac{1}{2}\vw^{\!\top} H \vw\), com
  \(H\) a \textbf{Hessiana} (curvaturas), autovalores \(\lambda_1 \ge \dots \ge
  \lambda_d > 0\): em cada autodireção, a quadrática 1D com \(a = \lambda_i\).

  \begin{teobox}{Convergência no caso multivariado}
    Converge sse \(\mdestaque{icMagenta}{\taxa < 2/\lambda_{\max}}\): a direção
    \emph{mais curva} dita o teto da taxa. A velocidade é limitada pela mais
    \emph{rasa}: com a melhor \(\taxa\) fixa, o erro cai como
    \(\bigl(\tfrac{\kappa - 1}{\kappa + 1}\bigr)^{t}\).
  \end{teobox}

  \begin{defbox}{Número de condição}
    \(\kappa = \lambda_{\max} / \lambda_{\min}\). Mede a \emph{anisotropia} do
    vale: \(\kappa = 1\) é uma tigela redonda; \(\kappa \gg 1\) é um desfiladeiro
    estreito e comprido.
  \end{defbox}

  \begin{itemize}\setlength{\itemsep}{0.2ex}
    \item Redes mal inicializadas ou sem normalização têm \(\kappa\) enorme —
          normalização de entrada, \emph{batch norm} e boas inicializações
          aceleram o treino porque \emph{reduzem} \(\kappa\)
  \end{itemize}
\end{frame}

% ------------------------------------------------------------
\begin{frame}{Por que $\kappa$ grande dói: o ziguezague}
  \begin{center}
    \scalebox{0.9}{%
    \begin{tikzpicture}
      % curvas de nível de um vale anisotrópico
      \foreach \s in {0.4, 0.8, 1.2, 1.6, 2.0}
        \draw[icAzulClaro!60] (0,0) ellipse ({2.2*\s} and {0.45*\s});
      \fill[icVerde] (0,0) circle (2.5pt) node[below=2pt] {\scriptsize mínimo};
      % trajetória em ziguezague
      \draw[->, very thick, icCarmim]
        (-4.0, 0.72) -- (-3.1,-0.62) -- (-2.4, 0.53) -- (-1.85,-0.45)
        -- (-1.42, 0.38) -- (-1.1,-0.32) -- (-0.85, 0.27) -- (-0.65,-0.22);
      \node[icCarmim, align=center] at (-3.2, 1.25)
        {\scriptsize gradiente puro:\\ \scriptsize oscila na direção curva};
      % trajetória com momento
      \draw[->, very thick, icMagenta, dashed]
        (-4.0,-0.9) .. controls (-2.6,-0.55) and (-1.2,-0.25) .. (-0.15,-0.05);
      \node[icMagenta] at (-2.2,-1.25) {\scriptsize com momento: amortece e acelera};
    \end{tikzpicture}}
  \end{center}

  \begin{itemize}\setlength{\itemsep}{0.3ex}
    \item O gradiente aponta \emph{perpendicular às curvas de nível}, não para o
          mínimo: na direção curva ele é gigante, na rasa é minúsculo
    \item A \(\taxa\) precisa caber na direção curva
          (\(\taxa < 2/\lambda_{\max}\)) — e aí o avanço na direção rasa fica
          lento demais
    \item Duas saídas: \textbf{(i)} usar uma taxa \emph{por direção}
          (Newton, Adam) ou \textbf{(ii)} acumular velocidade (\textbf{momento})
  \end{itemize}
\end{frame}

% ------------------------------------------------------------
\begin{frame}{Exercício}
  \begin{exbox}{1 — Regimes de convergência}
    Considere \(\erro(w) = 2w^2\) (logo \(\erro'(w) = 4w\)) com \(w_0 = 1\).
    \begin{itemize}\setlength{\itemsep}{0.2ex}
      \item[(a)] Com \(\taxa = 0{,}1\), calcule \(w_1, w_2, w_3\). O
                 comportamento é monotônico ou oscilante?
      \item[(b)] Com \(\taxa = 0{,}4\), calcule \(w_1, w_2, w_3\). O que mudou?
      \item[(c)] Com \(\taxa = 0{,}6\), calcule \(w_1, w_2, w_3\). O que acontece
                 com \(\erro(w_t)\)?
      \item[(d)] Qual \(\taxa\) leva ao mínimo em um único passo?
    \end{itemize}
  \end{exbox}

  \begin{ideiabox}
    Antes de calcular, preveja: aqui a curvatura é \(a = 4\), então o teto é
    \(\taxa < 2/4 = 0{,}5\) e o passo perfeito é \(\taxa = 1/4\).
  \end{ideiabox}
\end{frame}

% ============================================================
\section{Taxa de aprendizado na prática}

\begin{frame}{Diagnóstico pela curva de perda}
  \begin{center}
    \scalebox{0.85}{%
    \begin{tikzpicture}
      \draw[->, icCinza] (0,0) -- (6.4,0) node[below] {\scriptsize iterações};
      \draw[->, icCinza] (0,0) -- (0,3.0) node[left] {\scriptsize $\erro$};
      % muito alta: explode
      \draw[very thick, icCarmim]
        plot[domain=0:2.6, samples=40] (\x, {1.5 + 0.12*\x*\x + 0.15*sin(9*\x r)});
      \node[icCarmim, right] at (2.6, 2.5) {\scriptsize $\taxa$ alta demais: diverge};
      % alta: estagna oscilando
      \draw[very thick, icDourado]
        plot[domain=0:6, samples=80] (\x, {0.9 + 1.4*exp(-1.6*\x) + 0.09*sin(11*\x r)});
      \node[icDourado, right] at (4.6, 1.35) {\scriptsize alta: cai e estagna oscilando};
      % boa
      \draw[very thick, icVerde]
        plot[domain=0:6, samples=60] (\x, {0.25 + 2.05*exp(-0.9*\x)});
      \node[icVerde, right] at (4.6, 0.45) {\scriptsize boa};
      % baixa
      \draw[very thick, icAzulClaro]
        plot[domain=0:6, samples=60] (\x, {2.3 - 0.22*\x});
      \node[icAzulClaro, right] at (4.55, 1.0) {\scriptsize baixa demais: lenta};
    \end{tikzpicture}}
  \end{center}

  \begin{alertabox}
    A perda que \alert{oscila e cresce} não se resolve com mais épocas: é
    divergência geométrica (\(\taxa\) acima do teto \(2/\lambda_{\max}\)).
    Reduza a taxa — tipicamente por um fator de 3 a 10 — e recomece.
  \end{alertabox}
\end{frame}

% ------------------------------------------------------------
\begin{frame}{Taxa fixa é desperdício: \emph{schedules}}
  \framesubtitle{O padrão moderno: aquecimento linear + decaimento cosseno}

  \begin{defbox}{\emph{Schedule} de taxa de aprendizado}
    Uma função \(\taxa_t\) que varia ao longo do treino. O padrão dominante em
    redes modernas (incluindo \emph{transformers} e LLMs):
    {\small\[
      \taxa_t =
      \begin{cases}
        \taxa_{\max}\, \dfrac{t}{T_w} & t < T_w \quad \text{(aquecimento linear)}\\[1.2ex]
        \taxa_{\min} + \tfrac{1}{2}(\taxa_{\max}-\taxa_{\min})
          \Bigl(1 + \cos \pi \dfrac{t - T_w}{T - T_w}\Bigr) & t \ge T_w
          \quad \text{(cosseno)}
      \end{cases}
    \]}
  \end{defbox}

  \begin{itemize}\setlength{\itemsep}{0.3ex}
    \item \textbf{Por que aquecer?} No início a curvatura local é enorme e mal
          estimada; passos grandes destroem o treino
    \item \textbf{Por que decair?} No fim, taxa alta fica quicando nas paredes
          do vale sem afinar a solução
    \item Em resumo: \(\taxa_t\) grande explora, \(\taxa_t\) pequena refina
  \end{itemize}
\end{frame}

% ------------------------------------------------------------
\begin{frame}[fragile]{\emph{Schedule} em PyTorch}
  \begin{lstlisting}
opt = torch.optim.SGD(model.parameters(), lr=0.1, momentum=0.9)

def aquece_e_decai(t, T_w=500, T=10_000):       # fator multiplicativo sobre lr
    if t < T_w:
        return t / T_w                          # aquecimento linear
    p = (t - T_w) / (T - T_w)                   # progresso em [0, 1]
    return 0.5 * (1 + math.cos(math.pi * p))    # decaimento cosseno

sched = torch.optim.lr_scheduler.LambdaLR(opt, aquece_e_decai)

for x, y in loader:
    opt.zero_grad()          # PyTorch ACUMULA gradientes: sempre zere antes
    perda = criterio(model(x), y)
    perda.backward()         # retropropagacao (Aula 4)
    opt.step()               # w <- w - lr_t * grad
    sched.step()             # avanca o schedule
  \end{lstlisting}

  \smallskip
  Esquecer o \texttt{zero\_grad()} soma gradientes de lotes diferentes — um dos
  bugs mais comuns em treino.
\end{frame}

% ============================================================
\section{Momento}

\begin{frame}{A bola pesada}
  \framesubtitle{Polyak (1964): dê inércia ao ponto que desce a superfície}

  \begin{defbox}{Descida de gradiente com momento}
    Mantenha uma \emph{velocidade} \(\vet{v}\) que acumula gradientes passados
    com fator \(\beta \in [0,1)\):
    \[
      \vet{v}_{t+1} = \beta\, \vet{v}_t + \deriv{\erro}{\vw}\Bigl|_{\vw_t},
      \qquad
      \mdestaque{icMagenta}{\vw_{t+1} = \vw_t - \taxa\, \vet{v}_{t+1}}
    \]
    Valor típico: \(\beta = 0{,}9\) — cada gradiente ecoa por
    \(\approx 1/(1-\beta) = 10\) passos.
  \end{defbox}

  \begin{itemize}\setlength{\itemsep}{0.3ex}
    \item Na direção \textbf{rasa}, os gradientes apontam sempre para o mesmo
          lado: as contribuições se \emph{somam} e o passo efetivo cresce até
          \(\approx \taxa/(1-\beta)\)
    \item Na direção \textbf{curva}, os gradientes alternam de sinal: as
          contribuições se \emph{cancelam} e a oscilação é amortecida
    \item Resultado: o ziguezague vira uma trajetória suave vale abaixo
  \end{itemize}
\end{frame}

% ------------------------------------------------------------
\begin{frame}{Quanto o momento acelera — e a variante de Nesterov}
  \begin{teobox}{Aceleração em vales mal condicionados}
    Em quadráticas, o gradiente puro precisa de \(O(\kappa)\) iterações por
    fator de redução do erro; com momento bem ajustado, \(O(\sqrt{\kappa})\).
    Para \(\kappa = 10^4\): dez mil passos vs.\ cem.
  \end{teobox}

  \begin{defbox}{Momento de Nesterov (1983)}
    Avalie o gradiente no ponto \emph{para onde a inércia já está levando}:
    \[
      \vet{v}_{t+1} = \beta\, \vet{v}_t
        + \deriv{\erro}{\vw}\Bigl|_{\,\vw_t - \taxa\beta \vet{v}_t}
    \]
    O ``olhar adiante'' corrige a velocidade antes de exagerar.
  \end{defbox}

  \smallskip
  Momento \alert{não} é acessório: SGD com momento \(0{,}9\) + cosseno segue
  competitivo com o Adam em visão — muitas vezes generalizando melhor.
\end{frame}

% ------------------------------------------------------------
\begin{frame}[fragile]{Momento em cinco linhas (e no PyTorch)}
  Implementação manual — exatamente a convenção interna do PyTorch:

  \begin{lstlisting}
v = [torch.zeros_like(p) for p in model.parameters()]

with torch.no_grad():
    for p, vp in zip(model.parameters(), v):
        vp.mul_(beta).add_(p.grad)     # v <- beta*v + grad
        p.add_(vp, alpha=-lr)          # w <- w - lr*v
  \end{lstlisting}

  Equivalente pronto (e a versão de Nesterov):

  \begin{lstlisting}
opt = torch.optim.SGD(model.parameters(), lr=0.1,
                      momentum=0.9)                  # Polyak
opt = torch.optim.SGD(model.parameters(), lr=0.1,
                      momentum=0.9, nesterov=True)   # Nesterov
  \end{lstlisting}

  \smallskip
  {\small Alguns livros usam \(\vet{v} \leftarrow \beta\vet{v} - \taxa
  \nabla\erro\); \(\vw \leftarrow \vw + \vet{v}\). Mesma dinâmica,
  hiperparâmetros \emph{não} intercambiáveis — confira a convenção do artigo.}
\end{frame}

% ============================================================
\section{Panorama e curiosidades}

\begin{frame}{Para discutir em aula}
  \begin{quizbox}
    Durante o treino, a perda caiu bem por 2\,000 iterações e depois passou a
    oscilar sem cair, num patamar alto. Um colega sugere: ``deixa treinando a
    noite toda, com mais iterações ela desce''.

    \smallskip
    \opcoes{A sugestão resolve \;/\; não resolve — e por quê?}
  \end{quizbox}

  \pause

  \begin{respbox}
    \textbf{Não resolve.} Oscilar num patamar é o regime
    \(1/\lambda < \taxa < 2/\lambda\) das direções curvas: o ponto fica quicando
    entre as paredes do vale, e mais iterações só repetem o quique. O remédio é
    \alert{reduzir \(\taxa\)} (manualmente ou com um \emph{schedule} de
    decaimento) — e é exatamente por isso que o decaimento cosseno derruba a
    perda no fim do treino.
  \end{respbox}
\end{frame}

% ------------------------------------------------------------
\begin{frame}{Curiosidade histórica: os ancestrais dos métodos adaptativos}
  \begin{itemize}\setlength{\itemsep}{0.4ex}
    \item \textbf{Newton:} \(\vw \leftarrow \vw - H^{-1}\nabla\erro\), isto é,
          \(\taxa = 1/\text{curvatura}\) por direção. Perfeito em teoria;
          inviável com \(d\) na casa dos bilhões (\(H\) tem \(d^2\) entradas)
    \item \textbf{Quickprop (1988)} e \textbf{Rprop (1993):} passo \emph{por
          peso} ajustado por sinais de gradientes — avós diretos do RMSProp e
          do Adam, que veremos na próxima aula
    \item \textbf{Otimização sem gradiente} (algoritmos genéticos, estratégias
          evolutivas): dispensa retropropagação, mas escala mal com \(d\); hoje
          vive em nichos (busca de arquiteturas, recompensas não diferenciáveis).
          Voltaremos a ela só como panorama, no fim do curso
  \end{itemize}

  \begin{ideiabox}
    A linha mestra do deep learning moderno é clara: \textbf{tudo que puder ser
    diferenciável, será} — porque o gradiente é a única informação de busca que
    escala para trilhões de parâmetros.
  \end{ideiabox}
\end{frame}

% ------------------------------------------------------------
\begin{frame}{Resumo da aula}
  \begin{itemize}\setlength{\itemsep}{0.4ex}
    \item A superfície de perda de uma rede profunda é \textbf{não convexa}, mas
          em alta dimensão os vilões práticos são selas e platôs — não mínimos
          locais ruins
    \item Na quadrática, tudo se decide por \(\taxa\) vs.\ curvatura:
          converge sse \(\taxa < 2/\lambda_{\max}\); passo ótimo
          \(\taxa = 1/\lambda\)
    \item O \textbf{número de condição} \(\kappa\) mede a anisotropia do vale e
          dita a velocidade: \(O(\kappa)\) para gradiente puro,
          \(O(\sqrt{\kappa})\) com momento
    \item Na prática: diagnostique pela curva de perda; use
          \textbf{aquecimento + decaimento cosseno}; SGD com momento
          \(0{,}9\) é uma base forte até hoje
    \item Próxima aula: por que treinar com \emph{minilotes} (SGD), e os
          otimizadores adaptativos — RMSProp, Adam e AdamW
  \end{itemize}
\end{frame}

\end{document}
